\usepackage[a4paper,margin=20mm,headheight=16pt]{geometry}
\usepackage[T1]{fontenc}
\usepackage{lmodern}
\usepackage{microtype}
\usepackage{etoolbox}
\usepackage{atbegshi}
\usepackage[table]{xcolor}
\usepackage{graphicx}
\usepackage{array,booktabs,longtable,tabularx}
\usepackage{enumitem}
\usepackage{amsmath}
\usepackage{fancyhdr}
\usepackage{titlesec}
\usepackage{float}
\usepackage{pdflscape}
\usepackage{listings}
\usepackage{xurl}
\usepackage{hyperref}
\definecolor{navy}{HTML}{15324F}
\definecolor{teal}{HTML}{007F86}
\definecolor{mist}{HTML}{EEF5F7}
\definecolor{slate}{HTML}{46586B}
\definecolor{amber}{HTML}{9A5700}
\hypersetup{colorlinks=true,linkcolor=navy,urlcolor=teal,
  pdfauthor={ICT University ERP project team; AI-assisted draft},
  pdflang={en},bookmarksnumbered=true}
\urlstyle{tt}
\AtBeginShipout{\AtBeginShipoutAddToBox{\special{papersize=\the\paperwidth,\the\paperheight}}}
\setlength{\parindent}{0pt}
\setlength{\parskip}{5pt}
\setlength{\emergencystretch}{3em}
\setcounter{tocdepth}{1}
\renewcommand{\arraystretch}{1.18}
\setlist{nosep,leftmargin=*,topsep=4pt}
\titleformat{\section}{\Large\bfseries\color{navy}}{\thesection}{0.7em}{}
\titleformat{\subsection}{\large\bfseries\color{teal}}{\thesubsection}{0.7em}{}
\titleformat{\subsubsection}{\normalsize\bfseries\color{navy}}{\thesubsubsection}{0.7em}{}
\titlespacing*{\section}{0pt}{12pt}{8pt}
\titlespacing*{\subsection}{0pt}{10pt}{4pt}
\pagestyle{fancy}
\fancyhf{}
\fancyhead[L]{\small\color{slate}ICT UNIVERSITY \quad / \quad ERP}
\fancyhead[R]{\small\color{slate}\DocCode\ \textbar\ v1.0}
\fancyfoot[L]{\footnotesize\color{slate}SEN4121 \quad Summer 2026 \quad Review draft}
\fancyfoot[R]{\small\thepage}
\renewcommand{\headrulewidth}{0.5pt}
\newcolumntype{L}[1]{>{\raggedright\arraybackslash}p{#1}}
\newcolumntype{Y}{>{\raggedright\arraybackslash}X}
\newcommand{\code}[1]{\path{#1}}
\newcommand{\source}[1]{\par{\footnotesize\color{slate}\textbf{Evidence:} \path{#1}}\par}
\newcommand{\statusI}{\textbf{\textcolor{teal}{Implemented}}}
\newcommand{\statusP}{\textbf{\textcolor{amber}{Partial}}}
\newcommand{\statusT}{\textbf{\textcolor{slate}{Target}}}
\newcommand{\callout}[2]{\par\smallskip
  \noindent\fcolorbox{teal}{mist}{\parbox{\dimexpr\linewidth-2\fboxsep-2\fboxrule\relax}{
  \textbf{\color{navy}#1}\par #2}}\par\smallskip}
\newcommand{\cover}[2]{
  \hypersetup{pageanchor=false}
  \begin{titlepage}
  \vspace*{8mm}
  {\color{teal}\rule{\textwidth}{3pt}}\par\vspace{8mm}
  {\Large\bfseries\color{navy}ICT UNIVERSITY}\par
  {\large Faculty of Information and Communication Technologies}\par
  \vspace{18mm}
  {\Huge\bfseries\color{navy}#1\par}
  \vspace{6mm}{\Large\color{teal}Multi-tenant University ERP\par}
  \vspace{4mm}{\large Academic \textbullet\ Marketing \& Finance\par
  Administration \& Human Resources}\par
  \vspace{15mm}
  \begin{tabularx}{\textwidth}{@{}L{40mm}Y@{}}
  \toprule
  Document identifier & \DocCode\\
  Version / status & 1.0 / Evidence-grounded review draft\\
  Baseline & Backend \texttt{6e10c54}; frontend \texttt{6cc6859}; reviewed 13 September 2026\\
  Course & SEN4121 -- Large System Environment, Summer 2026\\
  Instructor & Engr.\ Tanwi Nkiamboh\\
  Scope & React frontend and four-service backend, with deployment and integration boundaries\\
  Prepared with & AI assistance using Copilot SDK in VS Code; team review and approval pending\\
  \bottomrule
  \end{tabularx}
  \vfill
  \callout{Reading guide}{#2}
  {\small This document describes observed code, specified requirements and
  unverified targets separately. It is not a certification of production readiness,
  statutory payroll compliance or successful execution of application tests.}
  \end{titlepage}
}
\newcommand{\frontmatter}[1]{
  \pagenumbering{roman}
  \hypersetup{pageanchor=true}
  \phantomsection
  \section*{Document control}
  \addcontentsline{toc}{section}{Document control}
  \begin{tabularx}{\textwidth}{@{}L{29mm}L{33mm}Y@{}}
  \toprule Version & Date & Change\\\midrule
  1.0 & 13 September 2026 & Initial code-grounded specification with editable PlantUML sources and local PDF build.\\
  \bottomrule
  \end{tabularx}
  \subsection*{Review and approval}
  \begin{tabularx}{\textwidth}{@{}L{44mm}Y L{27mm}@{}}
  \toprule Responsibility & Name / signature & Date\\\midrule
  Product / academic representative & Pending team nomination & Pending\\
  Finance and payroll representative & Pending team nomination & Pending\\
  Technical lead / document reviewer & Pending team nomination & Pending\\
  Course submission approval & Pending group review & Pending\\
  \bottomrule
  \end{tabularx}
  \subsection*{Authority and interpretation}
  The supplied four-page SEN4121 project brief is the source of stakeholder
  obligations. Source code, migrations and tests determine what is implemented.
  Existing operational reports are historical evidence, not rerun results.
  The ISO/IEC/IEEE 29148 reference is used as an organizational alignment for
  requirements; no conformance assessment or reproduction of that standard is claimed.
  \subsection*{Academic integrity disclosure}
  This specification and its diagrams were generated with AI assistance from the
  project brief and repository evidence. Group members must review, correct and
  understand the material, disclose this assistance at the defense, and provide
  their own signed originality declaration and individual contribution statements.
  No names, signatures, approvals, measurements or contributions have been invented.
  \subsection*{Pagination policy}
  #1 Main narrative pages are deliberately bounded; extensive requirements,
  interface inventories and diagrams are appendices. The brief limits the
  combined SRS and technical report to 20 pages excluding appendices and diagrams.
  The package budgets six SRS narrative pages, eight SDD narrative pages and three
  front-matter pages per document: 20 combined non-appendix pages including covers,
  control pages and contents. Appendices are explicitly labeled and excluded from
  that total; this is not two separate 20-page allowances.
  \clearpage\tableofcontents\clearpage
}
\newcommand{\diagram}[4][portrait]{
  \clearpage
  \begingroup
  \def\sheetwidth{297mm}\def\sheetheight{420mm}
  \ifstrequal{#1}{wide}{\def\sheetwidth{420mm}\def\sheetheight{297mm}}{}
  \ifstrequal{#1}{large}{\def\sheetwidth{420mm}\def\sheetheight{594mm}}{}
  \ifstrequal{#1}{large-wide}{\def\sheetwidth{594mm}\def\sheetheight{420mm}}{}
  \newgeometry{layoutwidth=\sheetwidth,layoutheight=\sheetheight,margin=18mm}
  \setlength{\paperwidth}{\sheetwidth}
  \setlength{\paperheight}{\sheetheight}
  \setlength{\headwidth}{\textwidth}
  \subsection{#3}
  \begin{figure}[H]\centering
  \includegraphics[width=\textwidth,height=0.84\textheight,keepaspectratio]{diagrams/rendered/#2.png}
  \caption{#3. #4}\label{fig:#2}
  \end{figure}
  {\footnotesize Editable source: \code{diagrams/#2.puml}. This is a large-format
  diagram sheet; use original size or the companion SVG for full-detail review.
  Solid relationships denote
  implemented structure unless otherwise marked; dashed notes distinguish logical,
  external or target elements.}
  \clearpage
  \endgroup
  \restoregeometry
}
\lstset{basicstyle=\ttfamily\footnotesize,breaklines=true,
  frame=single,rulecolor=\color{teal},backgroundcolor=\color{mist},
  columns=fullflexible,keepspaces=true,showstringspaces=false}
